\subsection{The bracket}

Let \(r \in N_{K}\) and let X be a manifold of class \(C^{r+1}\) .

8.5.1. We apply the definitions of no. 8.4, taking for \(\tau\) the identity functor: \(\tau(\mathbf{V}) = \mathbf{V}\) for every Banach space V and \(\tau(u) = u\) for every isomorphism \(u\) of Banach spaces. We then have \(\tau(\mathrm{T}(\mathbf{X})) = \mathrm{T}(\mathbf{X})\) . If \(\xi\) and \(\eta\) are two vector fields of class \(\mathbf{C}^r\) on \(\mathbf{X}\) , we set

\[
[ \xi , \eta ] = \theta_ {\xi}. \eta ;\tag{1}
\]

it is a vector field of class \(C^{r-1}\) on X, called the bracket of \(\xi\) and \(\eta\) .

8.5.2. The map \((\xi, \eta) \mapsto [\xi, \eta]\) is K-bilinear and alternating. If \(f\) and \(g\) are functions of class \(\mathbf{C}^r\) on \(\mathbf{X}\) , with values in \(\mathbf{K}\) , we have:

\[
[ f \xi , g \eta ] = f g [ \xi , \eta ] + (f D _ {\xi} g) \eta - (g D _ {\eta} f) \xi .\tag{2}
\]

8.5.3. Let F be a Banach space. We have

\[
\mathbf {D} _ {[ \xi , \eta ]} = \mathbf {D} _ {\xi} \circ \mathbf {D} _ {\eta} - \mathbf {D} _ {\eta} \circ \mathbf {D} _ {\xi}\tag{3}
\]

in the space of maps from \(\mathcal{C}^{r+1}(X;F)\) to \(\mathcal{C}^{r-1}(X;F)\) .

If \(r \geqslant 2\) and if \(\zeta\) is a third vector field of class \(C^r\) on X, we have

\[
[ [ \xi , \eta ], \zeta ] = [ \xi , [ \eta , \zeta ] ] - [ \eta , [ \xi , \zeta ] ] ^ {(1)}\tag{4}
\]

or also

\[
[ \xi , [ \eta , \zeta ] ] + [ \eta , [ \zeta , \xi ] ] + [ \zeta , [ \xi , \eta ] ] = 0.
\]

More generally, if \(r \geqslant 2\) and if \(\tau\) is a vector functor for isomorphisms, we have

\[
\theta_ {[ \xi , \eta ]} = \theta_ {\xi} \circ \theta _ {\eta} - \theta _ {\eta} \circ \theta_ {\xi}\tag{5}
\]

in the space of maps from \(\mathcal{C}^{r+1}(\mathrm{X};\mathrm{F})\) to \(\mathcal{C}^{r-1}(\mathrm{X};\mathrm{F})\) (with \(\mathrm{F} = \tau (\mathrm{T}(\mathrm{X}))\) ).

If \(r \geqslant \infty\) the vector fields of class \(C^r\) on X form a K-Lie algebra for the bracket (LIE, I, § 1, no. 2).

8.5.4. Let E be a Banach space, U an open subset of E, \(\xi\) and \(\eta\) two vector fields of class \(C^{r}\) on U, identified with elements of \(\mathcal{C}^{r}(U;E)\) . Their bracket is given by the formula:

\[
[ \xi , \eta ] = \mathbf {D} _ {\xi} \eta - \mathbf {D} _ {\eta} \xi .\tag{6}
\]

8.5.5. Let \((u^1, \ldots, u^n)\) be a coordinate system of X in an open set U. Let \(\xi = \sum a^i \frac{\partial}{\partial u^i}\) and \(\eta = \sum b^i \frac{\partial}{\partial u^i}\) be two vector fields of class \(C^r\) on U. Set \([\xi, \eta] = \sum c^i \frac{\partial}{\partial u^i}\) . One has

\[
c ^ {i} = \sum_ {1 \leqslant j \leqslant n} \left(a ^ {j} \frac {\partial b ^ {i}}{\partial u ^ {j}} - b ^ {j} \frac {\partial a ^ {i}}{\partial u ^ {j}}\right).\tag{7}
\]

8.5.6. Let \(\varphi: X \to Y\) be a morphism of manifolds of class \(C^{r+1}\) , \(\xi_1\) and \(\xi_2\) two vector fields on X, \(\eta_{1}\) and \(\eta_{2}\) two vector fields on Y, all of class \(C^r\) . If \(\xi_{1}\) and \(\xi_{2}\) are \(\varphi\)-related to \(\eta_{1}\) and \(\eta_{2}\) respectively, \([\xi_{1}, \xi_{2}]\) is \(\varphi\)-related to \([\eta_{1}, \eta_{2}]\) .

8.5.7. Let \(\omega\) be a differential form of degree \(p\) on X, with values in a Banach space F and of class \(\mathbf{C}^r\) , and let \(\xi_0, \ldots, \xi_p\) be vector fields of class \(\mathbf{C}^r\) on X. We have

\[
(\theta_ {\xi_ {0}}. \omega) (\xi_ {1}, \dots , \xi_ {p}) = D _ {\xi_ {0}} \omega (\xi_ {1}, \dots , \xi_ {p}) - \sum_ {i = 1} ^ {p} \omega (\xi_ {1}, \dots , [ \xi_ {0}, \xi_ {i} ], \dots , \xi_ {p})\tag{8}
\]

and

\[
\begin{array}{l} d \omega (\xi_ {0}, \dots , \xi_ {p}) = \sum_ {i = 0} ^ {p} (- 1) ^ {i} D _ {\xi_ {i}} \omega (\xi_ {0}, \dots , \hat {\xi} _ {i}, \dots , \xi_ {p}) \\ \qquad + \sum_ {i <   j} (- 1) ^ {i + j} \omega ([ \xi_ {i}, \xi_ {j} ], \xi_ {0}, \dots , \hat {\xi} _ {i}, \dots , \hat {\xi} _ {j}, \dots , \xi_ {p}). \end{array}\tag{9}
\]

If \(\xi\) and \(\eta\) are vector fields of class \(\mathbf{C}^r\) , one has

\[
\theta_ {\xi} \circ i (\eta) - i (\eta) \circ \theta_ {\xi} = i ([ \xi , \eta ])\tag{10}
\]

in the space of maps from \(^{r}\Omega^{p}(\mathbf{X};\mathbf{F})\) to \(^{r-1}\Omega^{p-1}(\mathbf{X};\mathbf{F})\) .

For \(p = 1\) , formula (9) becomes:

\[
d \omega (\xi , \eta) = D _ {\xi} \langle \eta , \omega \rangle - D _ {\eta} \langle \xi , \omega \rangle - \langle [ \xi , \eta ], \omega \rangle .\tag{11}
\]

